%% Elevark fra mattelab.no – laget automatisk av bygg.py. Rediger gjerne til eget bruk.
%% Kompiler med lualatex.
\documentclass[11pt,a4paper]{article}
\usepackage[norsk]{babel}
\usepackage[a4paper,margin=2.2cm,top=2cm,bottom=2.4cm]{geometry}
\usepackage{fontspec}
\usepackage{unicode-math}
\setmainfont{TeX Gyre Pagella}
\setsansfont{TeX Gyre Heros}[Scale=MatchLowercase]
\setmathfont{TeX Gyre Pagella Math}
\usepackage{xcolor}
\definecolor{blekk}{HTML}{16213A}
\definecolor{dempet}{HTML}{566077}
\definecolor{rute}{HTML}{CFC9B8}
\definecolor{m1}{HTML}{2F4AB0}
\definecolor{m2}{HTML}{99266D}
\usepackage[most]{tcolorbox}
\usepackage{qrcode}
\usepackage{needspace}
\usepackage{fancyhdr}
\usepackage[hidelinks]{hyperref}
\pdfvariable suppressoptionalinfo \numexpr 1+2+4+8+16+32+64+128+256+512\relax
\setlength{\parindent}{0pt}
\setlength{\parskip}{0.4em}
\pagestyle{fancy}
\fancyhf{}
\renewcommand{\headrulewidth}{0pt}
\fancyfoot[L]{\sffamily\footnotesize\color{dempet}mattelab.no · Rasjonale funksjoner og asymptoter · elevark uten fasit}
\fancyfoot[R]{\sffamily\footnotesize\color{dempet}Side \thepage}

\newcommand{\fase}[2]{\tcbox[on line, colframe=#2, colback=white, coltext=#2, boxrule=0.6pt, arc=2pt,
  boxsep=0pt, left=3pt, right=3pt, top=1.5pt, bottom=1.5pt, fontupper=\sffamily\bfseries\footnotesize]{#1}}
\newcommand{\steg}[4]{%
  \par\noindent{\color{rute}\rule{\linewidth}{0.4pt}}\par\nopagebreak
  \noindent\begin{minipage}[t]{2.1cm}\vspace{0pt}\fase{#1}{#2}\end{minipage}%
  \begin{minipage}[t]{\dimexpr\linewidth-2.1cm\relax}\vspace{0pt}#3\end{minipage}\par
  \vspace{#4}}
\newcommand{\oppgave}[2]{\par\Needspace{12\baselineskip}\vspace{0.6em}%
  {\fontsize{14}{17}\selectfont\bfseries{\color{m1}#1}\enspace #2}\par\vspace{0.2em}}

\begin{document}
\begin{minipage}[t]{0.72\linewidth}\vspace{0pt}
{\sffamily\small\color{dempet}Elevark · 1T og R1 · mattelab.no}\par\vspace{0.3em}
{\fontsize{24}{28}\selectfont\bfseries Rasjonale funksjoner og asymptoter}\par\vspace{1em}
{\sffamily\small Navn: \rule{4.6cm}{0.4pt}\quad Klasse: \rule{1.5cm}{0.4pt}\quad Dato: \rule{1.9cm}{0.4pt}}
\end{minipage}\hfill
\begin{minipage}[t]{0.22\linewidth}\vspace{0pt}\raggedleft
\qrcode[height=2.3cm]{https://mattelab.no/utforskere/rasjonale-funksjoner.html}\par\vspace{0.2em}
{\sffamily\scriptsize\color{dempet}Utforskeren}
\end{minipage}

\vspace{1em}
{\sffamily\small\color{dempet}Utforskeren ligger på \url{https://mattelab.no/utforskere/rasjonale-funksjoner.html}}\par
Hver oppgave har tre steg. \textbf{Gjett} først, uten å bruke utforskeren. \textbf{Test} så i utforskeren. \textbf{Forklar} til slutt med egne ord hvorfor det blir slik.

\oppgave{1}{Finn asymptotene}
\steg{Gjett}{m1}{Se på $f(x) = \frac{2x + 1}{x - 1}$. For hvilken $x$ er $f$ ikke definert? Hva tror du $f(x)$ er omtrent lik når $x = 100$?}{2.2cm}
\steg{Test}{m2}{Sett $a = 2$, $b = 1$, $c = 1$ og $d = -1$. Les av asymptotene. Prøv også $f(x) = \frac{1}{x}$.}{2.6cm}
\steg{Forklar}{blekk}{Forklar hvorfor den loddrette asymptoten er $x = -\frac{d}{c}$ og den vannrette er $y = \frac{a}{c}$.}{3.4cm}

\oppgave{2}{Nær den loddrette asymptoten}
\steg{Gjett}{m1}{Fortsatt med $f(x) = \frac{2x + 1}{x - 1}$: Hva tror du $f(1{,}01)$ og $f(0{,}99)$ er omtrent lik?}{2.2cm}
\steg{Test}{m2}{Sett opp funksjonen. Se på tabellen «Nær den loddrette asymptoten».}{2.6cm}
\steg{Forklar}{blekk}{Hvorfor blir verdiene så store? Hvorfor har de ulikt fortegn på hver side av $x = 1$?}{3.4cm}

\oppgave{3}{Et hull i grafen}
\steg{Gjett}{m1}{Se på $f(x) = \frac{2x - 4}{x - 2}$. Hva tror du grafen ser ut som? Har den en loddrett asymptote?}{2.2cm}
\steg{Test}{m2}{Start med $f(x) = \frac{2x - 3}{x - 2}$. Dra $b$ sakte fra $-3$ til $-4$. Følg med på grafen og på tallet $ad - bc$.}{2.6cm}
\steg{Forklar}{blekk}{Forkort brøken $\frac{2x - 4}{x - 2}$. Hvorfor er $f(2)$ ikke definert, selv om grafen ser ut som en rett linje? Hva er $\lim\limits_{x \to 2} f(x)$?}{3.4cm}

\oppgave{4}{Pris per t-skjorte}
\steg{Gjett}{m1}{Et trykkeri tar $200$ kr i oppstartspris og $15$ kr per t-skjorte. Gjennomsnittsprisen per skjorte for $x$ skjorter er $G(x) = \frac{15x + 200}{x}$. Hva tror du skjer med $G(x)$ når du bestiller mange skjorter? Kan prisen per skjorte bli under $15$ kr?}{2.2cm}
\steg{Test}{m2}{Sett $a = 15$, $b = 200$, $c = 1$ og $d = 0$. Zoom ut og se på grafen for $x > 0$. Finn $G(10)$ og $G(100)$.}{2.6cm}
\steg{Forklar}{blekk}{Hva betyr den vannrette asymptoten i praksis? Hvor mange skjorter må du bestille for at prisen per skjorte skal bli under $20$ kr? Hvilken del av grafen gir mening i denne situasjonen?}{3.4cm}

\oppgave{5}{Kan grafen krysse den vannrette asymptoten?}
\steg{Gjett}{m1}{Tror du grafen til $f(x) = \frac{ax + b}{cx + d}$ kan krysse den vannrette asymptoten $y = \frac{a}{c}$?}{2.2cm}
\steg{Test}{m2}{Start med standardoppsettet. Prøv mange ulike verdier av $a$, $b$, $c$ og $d$. Finner du en graf som krysser den vannrette asymptoten?}{2.6cm}
\steg{Forklar}{blekk}{Løs likningen $\frac{ax + b}{cx + d} = \frac{a}{c}$. Hva finner du ut?}{3.4cm}

\end{document}
